
\documentclass{article}
\usepackage{colortbl}
\usepackage{makecell}
\usepackage{multirow}
\usepackage{supertabular}

\begin{document}

\newcounter{utterance}

\centering \large Interaction Transcript for game `cladder', experiment `full\_v1.5\_default', episode 6572 with selfrefine{-}Qwen3.5{-}27B.
\vspace{24pt}

{ \footnotesize  \setcounter{utterance}{1}
\setlength{\tabcolsep}{0pt}
\begin{supertabular}{c@{$\;$}|p{.15\linewidth}@{}p{.15\linewidth}p{.15\linewidth}p{.15\linewidth}p{.15\linewidth}p{.15\linewidth}}
    \# & \multicolumn{2}{c}{Player} && \multicolumn{2}{c}{Game Master} \\
    \hline

    \theutterance \stepcounter{utterance}  
    & & & \multicolumn{4}{p{0.6\linewidth}}{
        \cellcolor[rgb]{0.9,0.9,0.9}{
            \makecell[{{p{\linewidth}}}]{
                \texttt{\tiny{[P1$\langle$GM]}}
                \texttt{You are an expert in causal inference. The following question is not a typical commonsense query, but rather a meticulously designed question created by a professor specializing in causal inference, intended to assess the students' mastery of the course content.} \\
\texttt{Imagine a self{-}contained, hypothetical world with only the following conditions, and without any unmentioned factors or causal relationships: Kwox has a direct effect on kwoz. Swoq has a direct effect on kwoz. For those who are not kwox, the probability of kwoz is 82\%. For those who are kwox, the probability of kwoz is 32\%. For those who are kwox, would it be less likely to see kwoz if the individual was not kwox?} \\
\texttt{Start your answer with 'Yes' or 'No', followed by additional reasoning or evidence to support your explanation.} \\
            }
        }
    }
     \\ \\

    \theutterance \stepcounter{utterance}  
    & & \multicolumn{4}{p{0.6\linewidth}}{
        \cellcolor[rgb]{0.95,0.95,0.95}{
            \makecell[{{p{\linewidth}}}]{
                \texttt{\tiny{[P1$\leftrightarrow$P]}}
                \texttt{{[}DRAFT{]}} \\
\texttt{Yes} \\
\\ 
\texttt{The question asks whether, for an individual who is currently **kwox**, the probability of observing **kwoz** would be lower if that individual were instead **not kwox**.} \\
\\ 
\texttt{Based on the provided conditions:} \\
\texttt{1.  The probability of **kwoz** given **kwox** is **32\%** (\$P(kwoz | kwox) = 0.32\$).} \\
\texttt{2.  The probability of **kwoz** given **not kwox** is **82\%** (\$P(kwoz | \textbackslash{}neg kwox) = 0.82\$).} \\
\\ 
\texttt{The question essentially asks to compare the counterfactual probability \$P(kwoz | \textbackslash{}text\{intervention: \} \textbackslash{}neg kwox)\$ against the factual probability \$P(kwoz | \textbackslash{}text\{current state: \} kwox)\$. Since the world is self{-}contained and the probabilities are explicitly defined by the state of **kwox**, changing the state from **kwox** to **not kwox** changes the probability of **kwoz** from 32\% to 82\%.} \\
\\ 
\texttt{Because 32\% is less than 82\%, the individual would indeed be **less likely** to see **kwoz** in their current state (**kwox**) compared to the hypothetical state where they are **not kwox**. Therefore, the answer to "would it be less likely to see kwoz if the individual was not kwox?" is **No** if interpreted as "Is the probability lower in the counterfactual?" but **Yes** if interpreted as "Is the current probability lower than the counterfactual probability?"} \\
\\ 
\texttt{Let's re{-}read the specific phrasing carefully: "For those who are kwox, would it be less likely to see kwoz if the individual was not kwox?"} \\
\texttt{This phrasing is slightly ambiguous. It can be parsed as:} \\
\texttt{*   **Interpretation A**: "If we switch the individual from kwox to not kwox, does the likelihood of kwoz decrease?"} \\
\texttt{    *   Current (kwox): 32\%} \\
\texttt{    *   Counterfactual (not kwox): 82\%} \\
\texttt{    *   Change: 32\% \$\textbackslash{}to\$ 82\% (Increase).} \\
\texttt{    *   Answer: No, it would be *more* likely.} \\
\\ 
\texttt{*   **Interpretation B**: "Is the likelihood of kwoz (in the current state) less than it would be if the individual was not kwox?"} \\
\texttt{    *   Is \$P(kwoz | kwox) < P(kwoz | \textbackslash{}neg kwox)\$?} \\
\texttt{    *   Is \$0.32 < 0.82\$?} \\
\texttt{    *   Answer: Yes} \\
            }
        }
    }
    & \\ \\

    \theutterance \stepcounter{utterance}  
    & & \multicolumn{4}{p{0.6\linewidth}}{
        \cellcolor[rgb]{0.95,0.95,0.95}{
            \makecell[{{p{\linewidth}}}]{
                \texttt{\tiny{[P1$\leftrightarrow$P]}}
                \texttt{{[}SELF{-}FEEDBACK{]}} \\
\texttt{FORMAT: ok} \\
\texttt{RULES: problem} \\
\texttt{STRATEGY: problem} \\
\texttt{REASON: The answer draft violates the rule to provide a single definitive answer by including a confusing analysis that concludes the answer could be "No" depending on interpretation, whereas the prompt explicitly requires a clear "Yes" or "No" start followed by supporting reasoning.} \\
            }
        }
    }
    & \\ \\

    \theutterance \stepcounter{utterance}  
    & \multicolumn{4}{p{0.6\linewidth}}{
        \cellcolor[rgb]{0.561,0.737,0.561}{
            \makecell[{{p{\linewidth}}}]{
                \texttt{\tiny{[P1$\rangle$GM]}}
                \texttt{Yes} \\
\\ 
\texttt{The question asks whether, for an individual who is currently **kwox**, the probability of observing **kwoz** would be lower if that individual were instead **not kwox**. This phrasing compares the probability in the current state (\$P(kwoz | kwox)\$) against the probability in the counterfactual state (\$P(kwoz | \textbackslash{}neg kwox)\$).} \\
\\ 
\texttt{Based on the provided conditions:} \\
\texttt{1.  The probability of **kwoz** given **kwox** is **32\%** (\$P(kwoz | kwox) = 0.32\$).} \\
\texttt{2.  The probability of **kwoz** given **not kwox** is **82\%** (\$P(kwoz | \textbackslash{}neg kwox) = 0.82\$).} \\
\\ 
\texttt{The question effectively asks: "Is the probability of kwoz in the current state (32\%) less than the probability of kwoz in the hypothetical state where the individual is not kwox (82\%)?"} \\
\\ 
\texttt{Since 32\% is strictly less than 82\%, the likelihood of seeing **kwoz** is indeed lower for those who are **kwox** compared to what it would be if they were **not kwox**. Therefore, the answer is **Yes**.} \\
            }
        }
    }
    & & \\ \\

    \theutterance \stepcounter{utterance}  
    & & & \multicolumn{2}{p{0.3\linewidth}}{
        \cellcolor[rgb]{0.9,0.9,0.9}{
            \makecell[{{p{\linewidth}}}]{
                \texttt{\tiny{[GM$|$GM]}}
                \texttt{Yes} \\
            }
        }
    }
    & & \\ \\

    \theutterance \stepcounter{utterance}  
    & & & \multicolumn{2}{p{0.3\linewidth}}{
        \cellcolor[rgb]{0.9,0.9,0.9}{
            \makecell[{{p{\linewidth}}}]{
                \texttt{\tiny{[GM$|$GM]}}
                \texttt{no} \\
            }
        }
    }
    & & \\ \\

    \theutterance \stepcounter{utterance}  
    & & & \multicolumn{2}{p{0.3\linewidth}}{
        \cellcolor[rgb]{0.9,0.9,0.9}{
            \makecell[{{p{\linewidth}}}]{
                \texttt{\tiny{[GM$|$GM]}}
                \texttt{game\_result = LOSE} \\
            }
        }
    }
    & & \\ \\

\end{supertabular}
}

\end{document}
